\scanpage{393}
\mspage{117}
% archive sheet number: 392
Je vais dégager à part ce qu'on vient d'obtenir :

\uline{Proposition} 7.7.1. Soit $(M,W)$ une catégorie de modèles, associée à un triple de Quillen clos \uline{spécial} $(W,C,F)$. Alors le prédérivateur \inserted{$D={}$} $D_{M,W}$ \inserted{sur $\mathrm{Cat}$} satisfait aux axiomes $\mathrm{Der}(1,2,3,4^*,5^*,6^*)$. Si de plus le triple de Quillen est \uline{très} spécial (donc en particulier $C=\mathrm{Mono}(M)$), alors il satisfait de plus à $\mathrm{Der}(4,5,6)$, i.e. c'est un dérivateur au plein sens du terme.

D'autre part, le localiseur $W=W_D$ dans $\mathrm{Cat}$ associé à $D$ (dans le cas d'un triple de Quillen spécial sans plus) \deleted{satisfait les} est un \uline{localiseur fondamental}, satisfaisant \deleted{l'axiome $\mathrm{Loc}(3\,\mathrm{bis}^{\deleted{\unsure{*|\prime}}})$.} \inserted{de plus} $\mathrm{Loc}(1\,\mathrm{ter},3\,\mathrm{bis})$.

\uline{NB} Ce qui est dit ici pour le cas d'un triple de Quillen très spécial (qui ne fait pas appel à la théorie des localiseurs
